\chapter{Fundamentals}\label{chap:fundamentals}

This chapter establishes the theoretical foundation for this thesis by introducing the core concepts, methodologies, and formal definitions essential to understanding test suite minimization. We present the metrics employed to quantify test suite effectiveness, provide a comprehensive description of the \emph{Stimulus-Response Matrix} (SRM) representation utilized within the LASSO platform, and formally define the test suite minimization problem within the context of computational complexity theory.

\section{Test Coverage and Mutation Score}

\emph{Code coverage} quantifies the proportion of program elements (statements, branches, paths, or conditions) exercised by a test suite~\cite{zhu1997}. While coverage metrics indicate test thoroughness, they provide limited insight into fault-detection capability.

\emph{Mutation testing} evaluates test suite quality by introducing small syntactic modifications (\emph{mutants}) into program source code~\cite{jia2011}. The \emph{mutation score} measures the percentage of mutants detected (killed) by the test suite:

\begin{equation}
\text{Mutation Score} = \frac{\text{Number of mutants killed}}{\text{Total number of mutants}} \times 100\%
\end{equation}

Mutation analysis provides superior insight into test effectiveness compared to coverage alone, as it directly measures fault-detection capability~\cite{fuzzingbook}.

\section{Stimulus-Response Matrices}\label{sec:srm}

The LASSO platform employs \emph{Stimulus-Response Matrices} (SRMs) as a unified, language-independent representation for capturing behavioral execution data across multiple software systems~\cite{Kessel2022}. An SRM is a two-dimensional matrix $M$ where rows represent stimuli (test sequences, method invocations, or API calls), columns represent executable software systems, and each matrix element $M_{ij}$ contains the structured response object produced by system $j$ when executed with stimulus $i$:

\begin{equation}
M_{ij} = \text{response object of system } j \text{ to stimulus } i
\end{equation}

Each response object $M_{ij}$ can contain rich, structured data including:
\begin{itemize}
  \item Textual output, return values, and exceptions
  \item Execution metrics (runtime, memory consumption)
  \item Code coverage data (statement, branch, path coverage percentages)
  \item Mutation testing results (mutation score, killed mutants)
  \item Execution traces and logs
\end{itemize}

\textbf{SRM Generation Mechanism:} SRMs are produced through LASSO's experimental execution pipeline using two key constructs~\cite{Kessel2022}:
\begin{enumerate}
  \item \emph{Sequence Sheets} define ordered execution steps, specifying how stimuli are constructed (e.g., method call sequences with parameters) and their dependencies
  \item \emph{Actuation Sheets} define the experimental setup, specifying which systems should be tested with which stimuli in LASSO's execution arena
\end{enumerate}

\subsubsection{Data Flow from Specification to Execution Results}

The relationship among these three artifacts follows a clear data flow pipeline:

\begin{enumerate}
  \item \textbf{Sequence Sheets (Stimulus Specification):} Each row defines one stimulus — an ordered sequence of actions, API calls, or method invocations with their parameters. Columns contain execution steps, input values, and expected outputs. Sequence Sheets function as the \emph{test case specification}, describing \emph{what} to execute. These specifications determine the row structure of the resulting SRM.

  \item \textbf{Actuation Sheets (Execution Schedule):} Each row defines one \emph{actuation} — a pairing of a system (software executable, container, or unit) with a stimulus from the Sequence Sheet. Actuation Sheets may include environment configuration parameters such as runtime limits, virtual machine images, or compiler versions. They function as the \emph{experiment execution schedule}, specifying \emph{where} and \emph{under what conditions} to execute each stimulus. These specifications determine the column structure of the resulting SRM.

  \item \textbf{Arena Execution:} LASSO's Arena component processes the Actuation Sheet, executing each (stimulus, system) pair according to the specifications in the Sequence Sheet and recording the resulting behavioral observations.

  \item \textbf{Stimulus-Response Matrices (Execution Results):} The Arena stores recorded responses as structured objects in SRM cells $M_{ij}$. Each cell captures the complete execution result for system $j$ responding to stimulus $i$, including outcome status, coverage metrics, mutation scores, runtime measurements, and exception logs.
\end{enumerate}

This pipeline can be summarized as: \emph{Sequence Sheet} (rows) + \emph{Actuation Sheet} (columns) $\xrightarrow{\text{Arena}}$ \emph{SRM} (cells with response objects). The SRM thus serves as a persistent, queryable repository of experimental results that can be analyzed for behavioral comparison, functional equivalence detection, or, as in this thesis, test suite optimization.

After execution, the resulting SRM records observed system behaviors at two levels of abstraction: (1) \emph{black-box view}, which captures only the final input/output parameters of sequence invocations, and (2) \emph{white-box view}, which records the complete method invocation trace including all intermediate inputs, outputs, and state changes. Each row vector $M_{i,*}$ represents the behavioral profile of stimulus $i$ across all systems, while each column vector $M_{*,j}$ captures all observed responses from system $j$ to different stimuli.

\textbf{Example:} Consider a simple SRM with 4 stimuli (test sequences) and 6 systems, where each response object contains execution results:
\begin{equation}
M = \begin{pmatrix}
\{\text{pass}, 85\%, 12\text{ms}\} & \{\text{pass}, 92\%, 8\text{ms}\} & \{\text{fail}, 78\%, 15\text{ms}\} & \cdots \\
\{\text{pass}, 90\%, 10\text{ms}\} & \{\text{pass}, 88\%, 9\text{ms}\} & \{\text{pass}, 95\%, 7\text{ms}\} & \cdots \\
\{\text{fail}, 65\%, 20\text{ms}\} & \{\text{pass}, 91\%, 11\text{ms}\} & \{\text{pass}, 87\%, 13\text{ms}\} & \cdots \\
\{\text{pass}, 88\%, 14\text{ms}\} & \{\text{fail}, 70\%, 18\text{ms}\} & \{\text{pass}, 93\%, 6\text{ms}\} & \cdots
\end{pmatrix}
\end{equation}

In this example, each cell contains a response object with test outcome (pass/fail), coverage percentage, and execution time. Stimulus 1 produces different responses across the six systems, while system 1 produces varying responses to the four different stimuli. Such behavioral profiles enable large-scale behavioral analysis including similarity measurements, code clone detection, functional equivalence analysis, and performance comparison across multiple systems. Figure~\ref{fig:srm-structure} illustrates the canonical SRM structure used for multi-system behavioral comparison.

\begin{figure}[H]
  \centering
  \begin{tabular}{c|cccccc}
    & \textbf{System 1} & \textbf{System 2} & \textbf{System 3} & \textbf{System 4} & \textbf{System 5} & \textbf{System 6} \\
    \hline
    \textbf{Stimulus 1} & \{pass, 85\%\} & \{pass, 92\%\} & \{fail, 78\%\} & \{pass, 89\%\} & \{pass, 91\%\} & \{pass, 87\%\} \\
    \textbf{Stimulus 2} & \{pass, 90\%\} & \{pass, 88\%\} & \{pass, 95\%\} & \{pass, 86\%\} & \{fail, 72\%\} & \{pass, 93\%\} \\
    \textbf{Stimulus 3} & \{fail, 65\%\} & \{pass, 91\%\} & \{pass, 87\%\} & \{pass, 94\%\} & \{pass, 89\%\} & \{fail, 70\%\} \\
    \textbf{Stimulus 4} & \{pass, 88\%\} & \{fail, 70\%\} & \{pass, 93\%\} & \{pass, 85\%\} & \{pass, 90\%\} & \{pass, 92\%\} \\
  \end{tabular}
  \caption{Canonical SRM structure: rows represent stimuli (test sequences or method invocations), columns represent different software systems, and cells contain structured response objects (e.g., test outcome and coverage percentage) enabling behavioral comparison across systems.}
  \label{fig:srm-structure}
\end{figure}

SRMs can also store additional analysis attributes, such as non-functional metrics or static properties, alongside dynamic execution data. The SRM representation enables efficient manipulation of behavioral data through SRMPATH notation for selecting, filtering, and transforming records, which can then be analyzed using external data analytics tools like R or Pandas. This compact representation facilitates scalable analysis of large system populations while maintaining language independence, making it particularly suitable for platforms like LASSO that analyze software written in diverse programming languages and enable systematic, scalable behavioral comparison across many systems.

\subsection{Leveraging SRMs for Test Suite Minimization}

Test suite minimization within LASSO operates directly on existing SRMs produced by the platform's execution pipeline~\cite{Kessel2022}. Rather than creating a new data structure, the minimization process extracts coverage-relevant information from SRM response objects for a selected target system. The workflow proceeds as follows:

\begin{enumerate}
  \item \textbf{System Selection:} Select one target system (column) from the SRM for minimization. In the current implementation, minimization focuses on a single system; multi-system minimization remains future work.
  
  \item \textbf{Coverage Extraction:} For each stimulus $i$ and the selected system $j$, extract coverage-related metrics from the response object $M_{ij}$. These metrics include:
  \begin{itemize}
    \item Statement, branch, or path coverage indicators
    \item Mutation testing results (killed mutants)
    \item Execution outcomes (pass/fail status)
  \end{itemize}
  
  \item \textbf{Algorithm Application:} Apply the HGS minimization algorithm (detailed in Chapter~\ref{chap:implementation}) using the extracted coverage data to select a subset of stimuli that preserve coverage requirements.
  
  \item \textbf{SRM Reduction:} Generate a reduced SRM containing only the selected stimuli (rows) while preserving the original SRM schema (systems as columns, structured response objects in cells). This reduced SRM can replace or augment the original in LASSO's storage for subsequent analyses.
\end{enumerate}

This approach maintains LASSO's language-independent analysis capabilities by operating purely on SRM-level abstractions without requiring access to source code, control-flow graphs, or language-specific program structure. Figure~\ref{fig:srm-adaptation} illustrates the minimization workflow showing system selection and stimulus reduction.

\begin{figure}[H]
  \centering
  \begin{minipage}[t]{0.45\textwidth}
    \centering
    \textbf{(a) Original SRM (Multi-System)}\\
    \vspace{0.3cm}
    \begin{tabular}{c|ccc}
      & \textbf{Sys1} & \textbf{Sys2} & \textbf{Sys3} \\
      \hline
      \textbf{Stim1} & \{pass, 85\%\} & \{fail, 70\%\} & \{pass, 90\%\} \\
      \textbf{Stim2} & \{pass, 92\%\} & \{pass, 88\%\} & \{pass, 95\%\} \\
      \textbf{Stim3} & \{fail, 65\%\} & \{pass, 91\%\} & \{fail, 72\%\} \\
      \textbf{Stim4} & \{pass, 88\%\} & \{pass, 87\%\} & \{pass, 93\%\} \\
      \textbf{Stim5} & \{pass, 90\%\} & \{fail, 68\%\} & \{pass, 89\%\} \\
    \end{tabular}
    \vspace{0.2cm}
    
    \small All stimuli, all systems
  \end{minipage}
  \hfill
  \begin{minipage}[t]{0.45\textwidth}
    \centering
    \textbf{(b) Reduced SRM (System 1)}\\
    \vspace{0.3cm}
    \begin{tabular}{c|ccc}
      & \textbf{Sys1} & \textbf{Sys2} & \textbf{Sys3} \\
      \hline
      \textbf{Stim1} & \{pass, 85\%\} & \{fail, 70\%\} & \{pass, 90\%\} \\
      \textbf{Stim2} & \{pass, 92\%\} & \{pass, 88\%\} & \{pass, 95\%\} \\
      \textbf{Stim4} & \{pass, 88\%\} & \{pass, 87\%\} & \{pass, 93\%\} \\
    \end{tabular}
    \vspace{0.2cm}
    
    \small Minimized stimuli, same schema
  \end{minipage}
  \caption{SRM-based test suite minimization workflow: (a) original SRM captures responses from all systems to all stimuli; (b) reduced SRM after minimization selects System 1 and eliminates redundant stimuli (Stim3, Stim5) while preserving SRM structure and coverage for the target system.}
  \label{fig:srm-adaptation}
\end{figure}

\section{The Test Suite Minimization Problem}

\subsection{Formal Problem Definition}
Let $T = \{t_1, t_2, \ldots, t_n\}$ denote a finite set of test cases (stimuli in LASSO terminology) and $R = \{r_1, r_2, \ldots, r_m\}$ represent the set of coverage requirements. Each test case $t_i \in T$ covers a subset of requirements $C(t_i) \subseteq R$. Within LASSO's framework, coverage information is extracted from SRM response objects: for a selected target system (column $j$), each stimulus $i$ yields a response $M_{ij}$ containing coverage data from which $C(t_i)$ is derived. The coverage achieved by the complete test suite is defined as:

\begin{equation}
C(T) = \bigcup_{i=1}^{n} C(t_i)
\end{equation}

The \emph{test suite minimization problem} seeks to find the minimum cardinality subset $S^* \subseteq T$ that preserves complete coverage:

\begin{equation}
S^* = \arg \min_{S \subseteq T} |S| \quad \text{subject to} \quad C(S) = C(T)
\end{equation}

\subsection{Computational Complexity}
This optimization problem is equivalent to the classical \emph{minimum set cover problem}, which is known to be NP-complete~\cite{yoo2012}. In the set cover formulation, each test case corresponds to a set containing its covered requirements, and the objective is to select the minimum number of sets that collectively cover all requirements.

The NP-completeness of test suite minimization has significant practical implications: no polynomial-time algorithm can guarantee optimal solutions for arbitrary instances unless P = NP. Consequently, practical approaches employ approximation algorithms~\cite{chvatal1979}, heuristic methods, or exact algorithms with exponential worst-case complexity that perform well on typical instances.

\subsection{Algorithmic Categories}
Existing approaches to test suite minimization can be taxonomically classified into several categories, each offering different trade-offs between solution quality, computational efficiency, and implementation complexity:

\begin{itemize}
  \item \textbf{Greedy heuristics:} These polynomial-time algorithms make locally optimal choices at each step. Examples include the classical greedy algorithm that iteratively selects tests covering the maximum number of uncovered requirements~\cite{chvatal1979}, and the Harrold--Gupta--Soffa (HGS) algorithm that prioritizes tests covering rare requirements~\cite{harrold1993}.
  
  \item \textbf{Exact optimization methods:} These approaches formulate minimization as integer linear programming (ILP) or constraint satisfaction problems, guaranteeing optimal solutions but with potentially exponential runtime complexity~\cite{mongiovi2020}.
  
  \item \textbf{Search-based and metaheuristic methods:} Evolutionary algorithms, genetic algorithms, and other population-based methods explore the solution space using fitness functions that balance multiple objectives such as suite size and coverage preservation~\cite{hussein2020}.
  
  \item \textbf{Data-driven and clustering techniques:} These methods exploit structural properties of extracted coverage information, using similarity measures, clustering algorithms, or machine learning techniques to identify representative test cases~\cite{bach2017}.
\end{itemize}

This taxonomic framework provides the organizational structure for the comprehensive literature review presented in Chapter~\ref{chap:literature-review}. Figure~\ref{fig:algorithm-taxonomy} illustrates the hierarchical classification of test suite minimization approaches.

\begin{figure}[H]
  \centering
  \begin{tabular}{|l|l|l|}
    \hline
    \textbf{Category} & \textbf{Approaches} & \textbf{Complexity} \\
    \hline
    \multirow{3}{*}{\textbf{Greedy Heuristics}} & Classical Greedy & $O(nm)$ \\
    & Harrold--Gupta--Soffa (HGS) & $O(nm \log m)$ \\
    & Delayed Greedy (Concept Analysis) & $O(n^2m)$ \\
    \hline
    \multirow{2}{*}{\textbf{Exact Optimization}} & Integer Linear Programming & Exponential \\
    & REDUNET (ILP + CFG) & Exponential \\
    \hline
    \multirow{2}{*}{\textbf{Search-Based}} & Genetic Algorithms & $O(gnm)$ \\
    & Quantum-Inspired GA & $O(gnm)$ \\
    \hline
    \multirow{2}{*}{\textbf{Data-Driven}} & Overlap-Aware Selection & $O(n^2m)$ \\
    & Clustering-Based & $O(n^2 + nm)$ \\
    \hline
  \end{tabular}
  \caption{Taxonomic classification of test suite minimization algorithms showing representative approaches and their computational complexity. Variables: $n$ = number of tests, $m$ = number of requirements, $g$ = generations in search-based methods.}
  \label{fig:algorithm-taxonomy}
\end{figure}
